
Le chapitre précédent construit simultanément les cinq opérations à partir du même état enrichi du continu. Chaque composante complète est déployée ici avant son application, avec fibre, extensions, relations, nombres, orientation, retenue, mémoire, frontière, parcours, multisection, survie, termes terminaux et lecteurs. La séparation expositive des cinq tuples n’en fait ni des branches, ni des facteurs, ni des objets indépendants : ce sont cinq caractéristiques intrinsèques de l’unique continu engendré, qui conservent littéralement le même état et le même registre. Le constructeur commun ne les réunit pas après coup ; il les produit simultanément comme structure interne d’un seul objet.

\section{Domaine, applications et compatibilités de l’objet terminal commun}

Pour exposer chaque caractéristique sans la séparer des autres, on utilise l’indice $T\in\{\mathrm{sp},\mathrm{sol},\mathrm{gau},\mathrm{coh},\mathrm{det}\}$, mais la donnée mathématique complète demeure le même continu enrichi. La vue interne n’est pas postulée de nouveau : c’est exactement la lecture typée construite dans \eqref{eq:pdf2-operacion-infinita-tipificada},
\begin{equation}\label{eq:pdf2-g-t-desde-continuo}
 \mathfrak G_T^{\rm int}:=
 \operatorname{View}_T\!\left(
 \mathcal O_\infty(\mathcal C_{\rm cont}^{\rm disc})\right).
\end{equation}
Sa présentation complète dans le canal \(T\) est
\begin{equation}\label{eq:pdf2-g-t-completo}
 \mathfrak G_T^{\rm int}=
 \bigl(
 \begin{gathered}
 \{\mathcal F_q^{(T)}\}_q,
 \{\operatorname{Ext}_\gamma^{(T)}\}_\gamma,
 \operatorname{Rel}_T,\operatorname{Num}_T,\operatorname{or}_T,\\
 \operatorname{Carr}_T,\operatorname{Mem}_T,\partial_T,\gamma_T,
 \mathfrak m_T,\operatorname{Surv}_T,X_{\chi_T},\tau_T,
 \operatorname{Term}_T,\operatorname{Lect}_T
 \end{gathered}
 \bigr).
\end{equation}
Il n’est pas permis de remplacer cet objet par $X_{\chi_T}$, par son cardinal ou par le nom d’un problème classique. Soit $\mathcal D_T^{\rm cl}$ l’espace des données classiques du problème correspondant. Celles-ci n’entrent que par
\begin{equation}\label{eq:pdf2-mapa-preparacion-aplicacion}
 \operatorname{App}_{T,d_T}:\mathfrak G_T^{\rm int}\longrightarrow
 \mathfrak G_T(d_T),\qquad d_T\in\mathcal D_T^{\rm cl},
\end{equation}
où $\mathfrak G_T(d_T)$ est, par définition, le couple formé par la vue interne complète et les coordonnées d’interface déployées ci-dessous. Ainsi, aucune fibre, relation, orientation, retenue, mémoire, frontière, parcours, multisection, survie ou terme terminal n’est perdu lors de l’introduction de la donnée classique. La projection interne $\operatorname{pr}_{X,T}:\mathfrak G_T^{\rm int}\twoheadrightarrow X_{\chi_T}$ est surjective. Si $\mathcal A_{T,d_T}:\mathfrak G_T(d_T)\to\mathcal M_T$, on définit
\[
 \widetilde{\mathcal A}_{T,d_T}:=
 \mathcal A_{T,d_T}\circ\operatorname{App}_{T,d_T},
 \qquad
 \mathfrak M_T:=\operatorname{Im}\mathcal A_{T,d_T}.
\]

\section{Déploiement interne des réalisations corrélatives avant leur identification aux problèmes classiques}

La formule générique précédente ne remplace pas l’écriture des cinq caractéristiques. Dans ce volume, chaque interface est matérialisée sans transformer la vue interne en objet ou domaine indépendant. Dans les cinq présentations suivantes, la première coordonnée est toujours $\mathfrak G_T^{\rm int}$ ; les listes de coordonnées spécifiques étendent donc le registre commun sans jamais s’y substituer. Nous utiliserons
\[
 d_{\rm sp}=(R,\mathcal D_R^{\rm cof}),\quad
 d_{\rm sol}=(u_0,f,T),\quad
 d_{\rm gau}=(\Lambda_m,g),\quad
 d_{\rm coh}=(X,p),\quad
 d_{\rm det}=E/\mathbb Q.
\]
Cette ligne fixe l’instant causal exact où les problèmes classiques apparaissent : après \(\mathfrak G_T^{\rm int}\), jamais au sein de sa génération HMT.

\subsection{Action spectrale–polaire du constructeur commun et conservation de la fibre critique}

La caractéristique spectrale--polaire conserve le registre
\begin{equation}\label{eq:pdf2-g-sp-explicito}
 \mathfrak G_{\rm sp}(d_{\rm sp})=
 \bigl(
 \begin{gathered}
  \mathfrak G_{\rm sp}^{\rm int};\\
  \widehat X_\infty,\Gamma_9,\widehat R=(R,K),
  \operatorname{Carr}_R,S_0,M,\mathcal O_9,\\
  \Xi_R,P_R,\mathcal D_R^{\rm cof},
  \mathcal F_R,\operatorname{Ext}_R,\operatorname{Surv}_R,
  \tau_R,\mathcal P_R
 \end{gathered}
 \bigr).
\end{equation}
Ici, $\mathcal D_R^{\rm cof}$ est le domaine des fonctions tests ; $\Xi_R$ exprime les canaux premier, gamma et polaire dans un même espace de Hilbert ; $P_R=P_R^*=P_R^2$ est la projection naturelle ; $S_0,M,\mathcal O_9$ conservent la mémoire et le terme terminal ; et $\widehat R=(R,K)$ conserve aussi bien la fenêtre que sa profondeur. L’objet de Weil est le complément, non un signe ajouté :
\begin{equation}\label{eq:pdf2-g-sp-complemento}
 B_{W,R}:=\Xi_R^*(I-P_R)\Xi_R.
\end{equation}
Les canaux régulés, la naturalité cofinale et le terme terminal appartiennent à la même caractéristique ; aucun n’est reconstruit à partir des zéros de $\zeta$. La structure interne fournit l’involution des feuillets, l’espérance des extensions, le défaut et le télescopage terminal. Le développement de référence exact précédent construit en outre l’interface de Weil : il fixe la contraction des feuillets dans \eqref{eq:amp-rh-owner-q}, relie le complément spectral--polaire dans \eqref{eq:amp-rh-owner-binding}, conserve le terme terminal au moyen de \eqref{eq:amp-rh-owner-telescopia} et prouve son extinction dans \eqref{eq:amp-rh-owner-terminal}. Dans la notation du constructeur commun, cette colligation déjà construite s’exprime comme
\begin{equation}\label{eq:pdf2-g-sp-coligacion-dato}
 \Sigma_R^{\rm sp}:
 \mathcal H_{+,R}\oplus\mathcal X_R^{\rm sp}
 \longrightarrow
 \mathcal H_{-,R}\oplus\mathcal L_R^{\rm sp}
 \oplus\mathcal X_R^{\rm sp},
 \qquad
 (\Sigma_R^{\rm sp})^*\Sigma_R^{\rm sp}=I,
\end{equation}
avec l’état terminal $x_{R,N}=q^NV_R^Nx_{R,0}$, $q=5/9$, et la sortie $K_R^{\rm sp}J_{+,R}=J_{-,R}$. Son identité d’énergie est
\begin{equation}\label{eq:pdf2-g-sp-identidad-finita-dato}
 \|J_{+,R}f\|^2
 =\|J_{-,R}f\|^2+
   \sum_{n=0}^{N-1}\|\ell_{R,n}f\|^2+
   q^{2N}\|x_{R,0}f\|^2.
\end{equation}
La factorisation \eqref{eq:amp-rh-owner-l} et son expression cofinale \eqref{eq:amp-rh-owner-publicacion} prouvent ces identités et l’entrelaceur premier--gamma--polaire avant cet assemblage. Le chapitre consacré à Riemann les déploie et les reconnaît au moyen du théorème~\ref{thm:amp-rh-cierre-cofinal-propietario} ; il ne construit pas a posteriori une interface absente et ne la reçoit pas des zéros de $\zeta$.

\subsection{Bilan solénoïdal du constructeur commun et conservation du flux}

Pour chaque coupure $M$ et profondeur $j$, l’image est
\begin{equation}\label{eq:pdf2-g-sol-explicito}
 \mathfrak G_{{\rm sol},M,j}(d_{\rm sol})=
 \bigl(
 \begin{gathered}
  \mathfrak G_{\rm sol}^{\rm int};\\
  \Xi_{M,j}^{\rm sol},C_{M\leftarrow N},M_{M,j}^+,M_{M,j}^-,
  D_{M,j},H_{M,j},E_9,J_9,P_9,Q_{\rm micro},\\
  \Gamma_{M,j}^{\rm term},\mathcal E_{M,j,T}^{\rm term},
  L_{M,j,T}^{\rm term},\ell_{M,j},\operatorname{Carr}_{M,j},\\
  \operatorname{Mem}_{M,j},\operatorname{Surv}_{M,j},\tau_{M,j}
 \end{gathered}
 \bigr).
\end{equation}
Le champ de Galerkin exprimé est une coordonnée de $\Xi_{M,j}^{\rm sol}$, mais l’état comprend en outre le défaut $C_{M\leftarrow N}$, les mémoires orientées, la counité et sa section, la dynamique terminale et la perte observée. Le Gram
\begin{equation}\label{eq:pdf2-g-sol-gram-estructural}
 U_{M,j;J,T}^{\rm term}
 :=(\mathbf G_{M,j;J,T}^{\rm term})^*
    \mathbf G_{M,j;J,T}^{\rm term}
\end{equation}
conserve dans $\mathbf G^{\rm term}$ la coordonnée terminale et toutes les pertes propagées. Le télescopage n’en efface aucune. Le développement de référence exact précédant cet assemblage construit l’application causale hydrodynamique complète : le relèvement PC--Fock de \eqref{eq:ns-import-pcf-lift}, la colligation isométrique de \eqref{eq:ns-import-cascade-isometry}, son télescopage de Stein \eqref{eq:ns-import-stein-telescope}, le registre \eqref{eq:ns-import-ledger}, la retenue comptabilisée une seule fois \eqref{eq:ns-import-carry-funded} et le retour exact \eqref{eq:ns-import-exact-return}. Dans ces mêmes coordonnées, les ondes de diffusion s’expriment comme
\begin{equation}\label{eq:pdf2-g-sol-ondas-dato}
 a_{M,j}=(2\sqrt\nu)^{-1}A_M^{-1/2}\Lambda^sf_{M,j},
 \qquad
 b_{M,j}=\sqrt\nu A_M^{1/2}\Lambda^su_{M,j}-a_{M,j},
\end{equation}
et la transition linéaire agit sur toutes les coordonnées mixtes $u^{\odot r}\otimes f^{\odot k}$. La racine commune et son transport ne sont pas postulés dans ce bloc : ils sont déjà construits par \eqref{eq:nsclose-owner-metric-definition}--\eqref{eq:nsclose-owner-root-uniformity}. Leur expression dans la carte commune est
\begin{equation}\label{eq:pdf2-g-sol-raiz-dato}
 J_T^{\rm sol}:\mathcal D_T^{\rm NS}\longrightarrow\mathcal H_T^{\rm sol},
 \qquad
 \mathbf G_{M,0;J,T}^{\rm term}\Lambda_{M,T}J_T^{\rm sol}
 =I_{M,J}J_T^{\rm sol},
 \qquad I_{M,J}^*I_{M,J}=I.
\end{equation}
La conservation conjointe du terme terminal, de la dissipation, de l’advection, de la retenue, du microétat, de la coquille et de la mémoire est prouvée par le télescopage de référence \eqref{eq:nsclose-owner-telescope} et son registre \eqref{eq:nsclose-proprietary-ledger}. Le budget uniforme, indépendant de $M,J$, est le théorème déjà établi par \eqref{eq:nsclose-proprietary-uniform-budget} et \eqref{eq:nsclose-uniform-sobolev-precompactness} ; il s’exprime ici dans la notation du constructeur commun comme
\begin{equation}\label{eq:pdf2-g-sol-cota-dato}
 \sup_{M,J}
 \|\mathbf G_{M,0;J,T}^{\rm term}\Lambda_{M,T}J_T^{\rm sol}d\|^2
 \le C_T^{\rm sol}\|d\|_{\mathcal D_T^{\rm NS}}^2,
\end{equation}
avec $C_T^{\rm sol}$ indépendant de $M,J$. Le chapitre ultérieur déploie la réalisation de Navier--Stokes et reconnaît son codomaine classique ; il ne fournit pas une prémisse absente ni ne sélectionne la borne à partir de la conclusion.

\subsection{Incidence projective de jauge du constructeur commun et descente au quotient}

L'image de jauge n'est pas la marginale des liaisons. Son objet complet par niveau est
\begin{equation}\label{eq:pdf2-g-gau-explicito}
 \mathfrak G_{{\rm gau},m}(d_{\rm gau})=
 \bigl(
 \begin{gathered}
  \mathfrak G_{\rm gau}^{\rm int};\\
  \Lambda_m,\widehat X_m,\widehat\mu_m^{\rm ov},
  \{\widehat\mu_{m,g}^{\rm YM},\nu_{m,g},K_{m,g}^{\rm phys}\}_{g>0},\\
  \mathcal G_m^{\rm full},\mathbb P_m^{\rm phys},
  \widehat J_m,\widehat\Gamma_m,
  \{E_{m,c}\}_c,\widehat H_m^{\rm ov},D_m^{\rm YM},\\
  \{\Theta_{\nu,m}\}_{\nu=1}^4,
  \{\mathfrak A_{+,\nu,m},\operatorname{ev}_{t}^{(\nu)},
      \mathcal V_{m,\nu}^{\rm OS}\}_{\nu=1}^4,\\
  \widehat\Omega_m^{(8)},P_{\Omega_m},W_c,
  \operatorname{Carr}_m,\operatorname{Mem}_m,
  \operatorname{Surv}_m,\tau_m,\mathcal P_m^{F^2}
 \end{gathered}
 \bigr).
\end{equation}
Le groupe $\mathcal G_m^{\rm full}$ agit simultanément sur les liens, les repères et les
coordonnées d’incidence $q_c\in\mathbb F_3^6$. La projection physique est sa moyenne
de Haar. Le générateur réduit est
\begin{equation}\label{eq:pdf2-g-gau-generador-estructural}
 D_m^{\rm YM}
 =\left.3\kappa_{\rm av}\sum_c(I-E_{m,c})
  \right|_{\operatorname{Ran}\mathbb P_m^{\rm phys}},
\end{equation}
et $W_c$ appartient à cette sous-algèbre. Les quatre réflexions agissent sur une même mesure, un même noyau et un même terminal.
Le propriétaire exact précédent construit l'action de jauge d'incidence \eqref{eq:amp-ym-accion-gauge-incidencia}, la projection physique de Haar \eqref{eq:amp-ym-proyeccion-fisica-haar} et la compression entrelacée \eqref{eq:amp-ym-compresion-entrelazada}. De cette construction proviennent le noyau complet équivariant $\widehat K_{m,t}^{\rm ov}$, sa descente au quotient physique $K_{m,t}^{\rm phys}$, les quatre opérateurs de réflexion et l'action $U_m(g)$ de $E(4)$ sur le cœur cylindrique lisse. Ce bloc assemble, sans les introduire comme données, le générateur, le vide, le témoin de courbure et les formes OS :
\begin{equation}\label{eq:pdf2-g-gau-entrelazamiento-dato}
 D_{m+1}^{\rm YM}\widehat J_m=\widehat J_mD_m^{\rm YM},
 \quad
 \mathbb P_{m+1}^{\rm phys}\widehat J_m
  =\widehat J_m\mathbb P_m^{\rm phys},
 \quad
 \widehat J_mW_{m,c}=W_{m+1,c}.
\end{equation}
La première et la troisième identités sont la publication, dans la notation
commune, de \eqref{eq:amp-ym-compresion-entrelazada} et de
\eqref{eq:amp-ym-testigo-fisico-gap} ; le semi-groupe, le vide simple et
l’écart exact s’obtiennent déjà dans
\eqref{eq:amp-ym-owner-semigrupo-vacio}--
\eqref{eq:amp-ym-owner-brecha-exacta}. Ici,
$\nu_{m,g}:=(\pi_m^{\rm phys})_*\widehat\mu_{m,g}^{\rm YM}$ et
$K_{m,g}^{\rm phys}(t)=e^{-tD_{m,g}^{\rm YM}}$ dans le quotient physique.
La complétude de jauge comprend en outre la désintégration selon chaque famille
d’hyperplans euclidiens. Si $A_j\in\mathfrak A_{+,\nu,m}$ et
$t_1\le\cdots\le t_n$, la marginale de la même mesure physique satisfait
\begin{equation}\label{eq:pdf2-g-gau-transfer-euclidean-dato}
\begin{aligned}
 &\int\prod_{j=1}^n
   (\tau_{t_je_\nu}A_j)(x)\,d\widehat\mu_{m,g}^{\rm YM}(x)\\
 &\quad=
 \int\prod_{j=1}^nA_j(y_j)\,
  \nu_{m,g}(dy_1)
  \prod_{j=1}^{n-1}K_{m,g}^{\rm phys}
       (t_{j+1}-t_j;y_j,dy_{j+1})\\
 &\quad=
 \left\langle\mathbf1,
  M_{A_1}e^{-(t_2-t_1)D_{m,g}^{\rm YM}}M_{A_2}\cdots
  e^{-(t_n-t_{n-1})D_{m,g}^{\rm YM}}M_{A_n}\mathbf1
 \right\rangle.
\end{aligned}
\end{equation}
La reconstruction précédente prouve l’identité dans les quatre directions
au moyen de \eqref{eq:amp-ym-cuatro-os} et construit un unique hamiltonien
entrelacé dans \eqref{eq:amp-ym-hamiltoniano-comun}. Par conséquent, la
translation OS est publiée ici comme
\begin{equation}\label{eq:pdf2-g-gau-os-generator-dato}
 \mathcal V_{m,\nu}^{\rm OS}T_{m,\nu}^{\rm OS}(s)
 (\mathcal V_{m,\nu}^{\rm OS})^{-1}
 =e^{-sD_{m,g}^{\rm YM}}.
\end{equation}
Ainsi, le générateur d’espérances et le Hamiltonien des translations ne sont pas
identifiés par leur nom : ce sont deux réalisations déjà construites de la même
désintégration, dont la forme limite et la publication de courbure figurent dans
\eqref{eq:amp-ym-forma-limite} et \eqref{eq:amp-ym-f2}.

\subsection{Cohérence cohomologique du constructeur commun et émergence de la classe globale}

Pour une variété projective lisse complexe $X$ et une codimension $p$, l’image cohomologique conserve
\begin{equation}\label{eq:pdf2-g-coh-explicito}
 \mathfrak G_{\rm coh}(d_{\rm coh})=
 \bigl(
 \begin{gathered}
  \mathfrak G_{\rm coh}^{\rm int};\\
  \{\mathscr S_N(X,p),\rho_{N+1,N}\}_N,
  \{\mathscr H_N^p(X),q_{N+1,N}\}_N,
  \{e_N\}_N,
  \{\iota_{N+1,N}\}_N,\\
  \{\mathcal G_{N+1}^{\rm coh},\mathcal R_{N+1}^{\rm coh},
       D_{N+1},a_{N+1},b_{N+1}\}_N,\\
  \mathcal F_{X,p},\operatorname{Ext}_{X,p},X_\chi(X,p),
  \operatorname{Rel}_{X,p},\operatorname{Carr}_{X,p},
  \operatorname{Mem}_{X,p},\partial_{X,p},\gamma_{X,p},\\
  \Sigma_{X,p},\operatorname{Surv}_{X,p},\tau_{X,p},
  R_{\rm Hdg},\operatorname{ch}^{\rm loc}_p,\operatorname{cl}_{X,p}
 \end{gathered}
 \bigr).
\end{equation}
Les carrés
\begin{equation}\label{eq:pdf2-g-coh-cuadrados}
 e_N\rho_{N+1,N}=q_{N+1,N}e_{N+1}
\end{equation}
relient observation, extension et raffinement. Le développement de référence exact précédent construit, sans la recevoir comme donnée, l’histoire basée avec ses corrections relatives, sa retenue, sa mémoire, sa frontière et son terme terminal. En particulier, \eqref{eq:amp-hodge-mapas-relativos}--\eqref{eq:amp-hodge-factorizacion-relativa} construisent les applications de l’incrément et \eqref{eq:amp-hodge-seccion-relativa} construit son inverse à droite. Dans la notation du constructeur commun, cette application relative s’exprime comme
\begin{equation}\label{eq:pdf2-g-coh-seccion-dato}
 s_{N+1}^{\rm rel}:\ker q_{N+1,N}\longrightarrow\ker\rho_{N+1,N},
 \qquad
 e_{N+1}^{\rm rel}s_{N+1}^{\rm rel}=I,
\end{equation}
naturelle sous raffinement par \eqref{eq:amp-hodge-naturalidad-seccion-relativa} et compatible avec la retenue et le terme terminal. La section initiale n’est pas non plus supposée : elle est construite dans \eqref{eq:amp-hodge-seccion-base}. Ce sont ces applications déjà démontrées, et non une classe finale choisie, que le bloc assemble et que le chapitre ultérieur reconnaît dans le codomaine de Hodge.

La complétude repose sur des actions et des identités construites dans le même développement de référence. La dégénérescence unitaire, la présentation relative et le lecteur sont assemblés ici au moyen de
\begin{equation}\label{eq:pdf2-g-coh-datos-operativos}
\begin{gathered}
 \iota_{N+1,N}:\mathscr S_N(X,p)\longrightarrow\mathscr S_{N+1}(X,p),
 \qquad \rho_{N+1,N}\iota_{N+1,N}=I,\\
 \mathscr A_{N+1}^{\rm rel}
  =\mathbb Q^{(\mathcal G_{N+1}^{\rm coh})}/\operatorname{im}D_{N+1},
 \qquad e_{N+1}^{\rm rel}a_{N+1}=b_{N+1},
 \qquad b_{N+1}D_{N+1}=0.
\end{gathered}
\end{equation}
L’ensemble $\mathcal G_{N+1}^{\rm coh}$ énumère exhaustivement tous les atomes nouveaux ainsi qu’une étiquette terminale par composante. La forme normale et l’inverse à droite sont construits par réduction basée dans \eqref{eq:amp-hodge-sigma-coordenada} ; leur compatibilité avec les raffinements est prouvée dans \eqref{eq:amp-hodge-naturalidad-seccion-relativa}. Ainsi, dans l’ordre selon profondeur, support, feuillet, parcours, retenue, mémoire et terme terminal, chaque atome d’observation possède une unique étiquette maximale et
\begin{equation}\label{eq:pdf2-g-coh-tabla-dato}
 b_{N+1}(c_\mu)=\bar c_\mu+
   \sum_{\nu<\mu}B_{\nu\mu}\bar c_\nu,
 \qquad B_{\nu\mu}\in\mathbb Q.
\end{equation}
La triangularité est une conséquence de cette construction basée, antérieure à toute classe $\alpha$ ; ce n’est ni une hypothèse ni un tableau fourni par le problème terminal.

La complétude et le lecteur du même objet sont typés par
\begin{equation}\label{eq:pdf2-g-coh-limite-lector-dato}
 X_\chi(X,p)=\varprojlim_N\mathscr S_N(X,p),
 \qquad
 q_N\operatorname{ev}^{\rm Hdg}_{\infty,X,p}
  =e_N\operatorname{pr}_N,
 \qquad
 \operatorname{cl}_{X,p}\operatorname{ch}^{\rm loc}_pR_{\rm Hdg}
  =\operatorname{ev}^{\rm Hdg}_{\infty,X,p}.
\end{equation}
Les projections $q_N$, y compris la coordonnée terminale, séparent les points du codomaine d’évaluation. L’extension opératoire \eqref{eq:amp-hodge-extension-operativa}--\eqref{eq:amp-hodge-ext}, le passage à la limite \eqref{eq:amp-hodge-xi-limite} et l’expression finale \eqref{eq:amp-hodge-publicacion-final} prouvent déjà que toute classe de Hodge arbitraire admet une histoire compatible. Le chapitre d’application déploie et reconnaît cette expression ; il ne fournit pas l’existence comme prémisse ultérieure.

\subsection{Droite déterminantale du constructeur commun et évaluation arithmétique terminale}

Pour une courbe elliptique $E/\mathbb Q$, l’image déterminantale conserve
\begin{equation}\label{eq:pdf2-g-det-explicito}
 \mathfrak G_{\rm det}(d_{\rm det})=
 \bigl(
 \begin{gathered}
  \mathfrak G_{\rm det}^{\rm int};\\
  \mathcal T_m^{\rm surv},G_{m,n},\Delta_{\rm AW},\varepsilon_{\rm AW},
  P_E^{\rm gen},1_E,P_E^K,P_E^{\rm PT},\\
  z_K,z_{\rm PT},X^{\rm orth},p_{\rm root},
  \mathcal C_E^K,\mathcal C_E^{\rm Iw},\mathcal C_E^{\rm PT},\\
  \ell_E,\mathcal C_E^{\rm loc,red},S_E(T),
  \operatorname{Carr}_E,\operatorname{Mem}_E,
  \operatorname{Surv}_E,\tau_E,\mathcal L_E
 \end{gathered}
 \bigr).
\end{equation}
L’unité $1_E$ est coexprimée dans les branches de Kato et de Poitou--Tate avant la dualité et la racine. Le produit fibré conserve l’image modulaire et l’histoire profonde. Le cône de localisation, les pages de Bockstein et la droite déterminante appartiennent à une même multisection basée ; ce ne sont pas des invariants choisis séparément après connaissance du rang ou du terme principal. Le développement de référence exact précédent construit la décomposition basée $\Theta_E$, l’homotopie du facteur fini--singulier, le comparateur complet Kato--FERS, les applications de pages et de recollement de Poitou--Tate et les applications de cocycles locales--globales. La chaîne commence dans la droite déterminante \eqref{eq:amp-bsd-linea-determinante}, construit l’application de chaînes \eqref{eq:amp-bsd-phi-kato-fers}, le comparateur de Poitou--Tate \eqref{eq:amp-bsd-comparador-pt} et le triangle local \eqref{eq:amp-bsd-triangulo-localizacion}. Ses identités opératoires s’expriment ici comme
\begin{equation}\label{eq:pdf2-g-det-identidades-dato}
\begin{gathered}
 \Phi_{K\leftrightarrow{\rm FERS}}\Psi_{K\leftrightarrow{\rm FERS}}=I,
 \qquad
 \Psi_{\rm PT}^{-1}\Psi_{\rm PT}=I,
 \qquad
 \Psi_{\rm PT}\Psi_{\rm PT}^{-1}=I,\\
 \ker\partial_{\rm loc}=\operatorname{im}\iota_{\Sha},
 \qquad
 \pi_\Phi\operatorname{Lift}_\Phi=I.
\end{gathered}
\end{equation}
Ces identités sont des conséquences de l’ensemble exact \eqref{eq:amp-bsd-paquete-identidades}--\eqref{eq:amp-bsd-paquete-composicion}. Le comparateur total est construit dans \eqref{eq:amp-bsd-comparador-total} et conserve la base commune du côté analytique jusqu’au côté arithmétique ; rang et terme principal sont exprimés ensuite dans \eqref{eq:amp-bsd-igualdad-grados}--\eqref{eq:amp-bsd-termino-principal}. Le chapitre BSD développe et reconnaît cette composition sans sélectionner rétrospectivement aucune de ses sorties.

\section{\(5\) lecteurs terminaux non permutables du continu généalogique}

L’application d’expression $\mathcal P_T$ est le lecteur mathématique qui traduit la structure dans le codomaine classique. La chaîne propre à ce volume est
\begin{equation}\label{eq:pdf2-cinco-cadenas}
 \mathfrak G_T^{\rm int}\xrightarrow{\operatorname{App}_{T,d_T}}
 \mathfrak G_T(d_T)\xrightarrow{\mathcal A_{T,d_T}}\mathfrak M_T
 \xrightarrow{\mathcal P_T}\mathcal C_T.
\end{equation}
Les correspondances sont fixées une par une :
\[
 \mathcal C_{\rm sp}=\mathrm{RH},\quad
 \mathcal C_{\rm sol}=\mathrm{NS},\quad
 \mathcal C_{\rm gau}=\mathrm{YM},\quad
 \mathcal C_{\rm coh}=\mathrm{Hodge},\quad
 \mathcal C_{\rm det}=\mathrm{BSD}.
\]
Dans l’assemblage actif, les développements de référence exacts qui précèdent cette section construisent également les cinq expressions finales : RH dans le théorème~\ref{thm:amp-rh-cierre-cofinal-propietario}, Navier--Stokes dans les théorèmes~\ref{thm:nsclose-proprietary-full} et \ref{thm:nsclose-global-smooth}, Yang--Mills dans le théorème~\ref{thm:amp-ym-terminacion}, Hodge dans le théorème~\ref{thm:amp-hodge-completitud-coinductiva} et BSD dans le théorème~\ref{thm:amp-bsd-comparadores-construidos}. Cette section enregistre leur assemblage corrélatif dans l’unique continu généalogique ; les chapitres ultérieurs déploient les cinq réalisations et leur reconnaissance classique, sans fournir de prémisses absentes ni utiliser les conclusions comme entrées.

\begin{falsifier}[Ablation de la donnée initiale]
Si une démonstration commence dans $X_{\chi_T}$, transforme une coordonnée de $\mathfrak G_T^{\rm int}$ en sélecteur facultatif ou remplace une fibre par son cardinal, elle n’a pas abrégé le même objet : elle l’a remplacé. La conclusion n’est pas transférable au substitut.
\end{falsifier}
